% ===================================================================
%  图表第 9 号 — 山与温泉小城。— 森林与打猎。
%  Traduction 2026 d'apres texto/io, controlee sur texto/fr, texto/en
%  et texto/es. Consigne : voir texto/zh/10-tubiao-01.tex.
% ===================================================================

%%K t09-apar-1 apar
\VUsaut{4.0mm}
\VUtitre{15.0pt}{60}{图表第 9 号}
\VUsaut{3.0mm}

%%K t09-tit-1 sub
\VUcentre{12.0pt}{0}{\textit{第一场。}}
\VUsaut{3.0mm}

%%K t09-tit-2 sub
\VUpk{11.4pt}{40}{山。--- 温泉小城。}
\VUsaut{3.0mm}

%%K t09-c1-01-1 p
1. --- 我到普拉茨一个星期了。我站在比利牛斯山那奇妙的\VUgras{山脉}
\textsuperscript{(1)} 面前时的惊叹，实在说不出来；那\VUgras{山脊} \textsuperscript{(2)}
在蓝天上刻出的轮廓，像一把巨大的锯。

%%K t09-c1-02-1 p
2. --- 我没想到比利牛斯山的\VUgras{山峰} \textsuperscript{(3)} 能有这样的
\VUgras{高度}，那些壮丽的\VUgras{冰川} \textsuperscript{(5)} 和积雪的
\VUgras{尖顶} \textsuperscript{(4)} 叫我看呆了，可怕的\VUgras{雪崩}就是从那上面
冲下来的。

%%K t09-tit-3 sub
\VUsaut{3.0mm}
\VUpk{10.2pt}{40}{山中的景色。}
\VUsaut{3.0mm}

%%K t09-c1-03-1 p
3. --- 到的第二天我就去了普拉茨附近那座早已熄了的\VUgras{火山}
\textsuperscript{(6)}。在\VUgras{火山口}光秃贫瘠的边上，还清清楚楚看得见几百年
前顺着\VUgras{山坡} \textsuperscript{(7)} 流下来的\VUgras{熔岩}留下的痕迹。在一处
\VUgras{斜坡}上，我居然找到了几块那种熔岩。

%%K t09-c1-04-1 p
4. --- 普拉茨在边界上，那座\VUgras{要塞} \textsuperscript{(8)} 守着法国和
西班牙之间的\VUgras{山口} \textsuperscript{(27)}，很像莱茵河边那些老城堡，它们仿佛在自己
的岩石顶上窥伺\VUgras{猎物}。

%%K t09-c1-05-1 p
5. --- 一只巨大的\VUgras{鹰} \textsuperscript{(9)} 常常引我注意，它的巢就在
\VUgras{堡} \textsuperscript{(8)} 不远的地方。这只\VUgras{猛禽}长着有力的
\VUgras{喙}，常常用\VUgras{爪}抓着一只羊羔或小山羊带给雏鸟。它的\VUgras{翅}
那么大，驮着重物还能滑翔很久。

%%K t09-tit-4 sub
\VUsaut{3.0mm}
\VUpk{10.2pt}{40}{步行的人。}
\VUsaut{3.0mm}

%%K t09-c1-06-1 p
6. --- 这里最惬意的一趟路是去科的\VUgras{瀑布} \textsuperscript{(10)}。
\VUgras{落下的水}打着旋儿泻下来，随后化作一条好看的\VUgras{溪}，消失在城
西的\VUgras{杉林} \textsuperscript{(11)} 里。我们穿过这片林子一直爬到
\VUgras{气象站} \textsuperscript{(12)}，从\VUgras{瞭望台}顶上把全区的
\VUgras{全景}尽收眼底。这\VUgras{景色}极好，\VUgras{大自然}像是把它手边所有
的宝贝都倾在这片地方了。

%%K t09-c1-07-1 p
7. --- 下山我们坐了\VUgras{缆车} \textsuperscript{(13)}，在一个小\VUgras{站}停下，
那站就在从前普拉茨领主住过的一座老\VUgras{封建城堡}的\VUgras{废墟}
\textsuperscript{(14)} 旁边。

%%K t09-c1-08-1 p
8. --- 可怜的城堡！它看见脚下那座整洁的\VUgras{温泉小城} \textsuperscript{(15)}
如今成了时髦的\VUgras{温泉胜地}，心里该怎么想呢？

%%K t09-c1-08-2 p
这里的\VUgras{气候}那么宜人，\VUgras{矿泉水}那么有益，在\VUgras{疗养院}
\textsuperscript{(17)} 里做的一\VUgras{疗程}简直是一桩乐事。

%%K t09-tit-5 sub
\VUsaut{3.0mm}
\VUpk{10.2pt}{40}{一座宜人的温泉小城。}
\VUsaut{3.0mm}

%%K t09-c1-09-1 p
9. --- 的确，为了让人住得舒服，什么都做到了。修了一座大\VUgras{高架桥}
\textsuperscript{(16)} 通铁路；\VUgras{矿泉厅}的\VUgras{园子} \textsuperscript{(18)}
布置得很雅致，那里还有\VUgras{音乐会}。几幢好看的\VUgras{别墅} \textsuperscript{(19)}
盖在\VUgras{赌场} \textsuperscript{(21)} 四周，一条养护得很好的\VUgras{公路}
\textsuperscript{(22)} 使来往极其方便。

%%K t09-c1-10-1 p
10. --- 一道简单的\VUgras{路堤} \textsuperscript{(23)} 把\VUgras{路} \textsuperscript{(20)}
同\VUgras{山谷} \textsuperscript{(25)} 本身隔开，谷底躺着一个秀气的小
\VUgras{湖} \textsuperscript{(26)}，湖水养着一条清亮的\VUgras{溪} \textsuperscript{(24)}。

%%K t09-c1-11-1 p
11. --- 山谷那边在开一座铁\VUgras{矿} \textsuperscript{(28)}。我好几次去看
\VUgras{矿工}下\VUgras{竖井}，还进过\VUgras{巷道}，他们整天在那里挖含着
\VUgras{金属}的\VUgras{矿石}。

%%K t09-c1-12-1 p
12. --- 矿旁边建了一座大\VUgras{冶炼厂}，好就地用上采出来的财富。
\VUgras{高炉} \textsuperscript{(29)} 用这里很多的\VUgras{煤}烧，能又快又便宜地
炼出\VUgras{生铁}。

%%K t09-c1-13-1 p
13. --- 离矿不远在开一座\VUgras{石场} \textsuperscript{(30)}。老\VUgras{采石工}
用\VUgras{镐}凿下来的既不是\VUgras{花岗岩}，也不是\VUgras{斑岩}，更不是
\VUgras{砂岩}，而是盖房子用的普通\VUgras{建筑石}。

%%K t09-c1-14-1 p
14. --- 这地方最好的装点之一是\VUgras{杉树} \textsuperscript{(31)}。哪里都有它
——\VUgras{峡谷} \textsuperscript{(32)} 底、\VUgras{急流} \textsuperscript{(33)} 边、
\VUgras{深壑} \textsuperscript{(34)} 或悬崖旁——它细细的针叶都添上一笔绿。

%%K t09-tit-6 sub
\VUsaut{3.0mm}
\VUpk{10.2pt}{40}{走路。}
\VUsaut{3.0mm}

%%K t09-c1-15-1 p
15. --- 我趁着假期多走路。盘山的\VUgras{小路} \textsuperscript{(35)} 使人不知不觉
就走了好几\VUgras{公里}，常常是好几\VUgras{法里}。

%%K t09-c1-15-2 p
\VUgras{里程碑} \textsuperscript{(36)} 和\VUgras{路标} \textsuperscript{(37)} 给小路
做着记号；路上常遇见一个年轻的\VUgras{山里人} \textsuperscript{(38)}，腰间挂着
\VUgras{行囊} \textsuperscript{(40)}，手里捧着一只\VUgras{旱獭} \textsuperscript{(39)}；
或者遇见一个\VUgras{吉卜赛人} \textsuperscript{(41)}，他的\VUgras{毛皮帽}
\textsuperscript{(42)} 同他那件破旧的\VUgras{斗篷} \textsuperscript{(43)} 很不相称。他
用\VUgras{链子}牵着一只会耍把戏的\VUgras{熊} \textsuperscript{(44)}。

%%K t09-tit-7 sub
\VUpk{10.2pt}{40}{登山。}
\VUsaut{3.0mm}

%%K t09-c1-16-1 p
16. --- 我差不多每天都跟\VUgras{向导} \textsuperscript{(48)} 去\VUgras{登山}；
背着\VUgras{背包} \textsuperscript{(47)}、手拿\VUgras{登山杖}，像个真正的
\VUgras{游客} \textsuperscript{(46)} 那样去攻那些\VUgras{岩石} \textsuperscript{(45)}
的时候，我很得意。

%%K t09-c1-17-1 p
17. --- 从山顶上看，平地上的人不比蚂蚁大；一群\VUgras{羊} \textsuperscript{(50)} 在\VUgras{牧场}
\textsuperscript{(49)} 上吃草，像小孩的玩具。一棵\VUgras{松树}
\textsuperscript{(51)}，甚至一座木头\VUgras{山屋} \textsuperscript{(52)}，也不过是绿地
上的一个小点。

%%K t09-c1-18-1 p
18. --- 走这些路时我们常在\VUgras{山道} \textsuperscript{(53)} 上遇见
\VUgras{猎羚人} \textsuperscript{(54)}，肩上扛着刚打到的野物。

%%K t09-c1-19-1 p
19. --- 我把一片很难得的\VUgras{蕨} \textsuperscript{(55)} 叶和一枝好看的粉色
\VUgras{石南} \textsuperscript{(57)} 放进了标本册，那是我上一次散步时在一个小
\VUgras{洞} \textsuperscript{(56)} 口采的。

%%K t09-tit-8 sub
\VUsaut{3.0mm}
\VUcentre{12.0pt}{0}{\textit{第二场。}}
\VUsaut{3.0mm}

%%K t09-tit-9 sub
\VUpk{11.4pt}{40}{森林。--- 打猎。}
\VUsaut{3.0mm}

%%K t09-c2-01-1 p
1. --- 昨天我在林子里过了极愉快的一天。住在那里的动物和\VUgras{昆虫}
\textsuperscript{(5)} 那忙碌的生活，叫我很感兴趣。

%%K t09-c2-01-2 p
\VUgras{蜘蛛} \textsuperscript{(1)} 在两根枝子中间结\VUgras{网} \textsuperscript{(2)}，
好捉飞过的\VUgras{苍蝇} \textsuperscript{(3)}；\VUgras{蜗牛} \textsuperscript{(4)} 驮着
壳吃力地爬；锹甲在四处找猎物；勤快的蚂蚁在\VUgras{蚁垤} \textsuperscript{(6)}
旁边使劲干活——这些小东西来来去去，忙得出奇。

%%K t09-c2-02-1 p
2. --- 一条\VUgras{蛇} \textsuperscript{(7)}，我起初当是\VUgras{毒蛇}，其实不过是
条不咬人的\VUgras{水蛇}，盘在一根树干底下，时不时伸出细细的小脑袋，看着
总像要咬人似的。我面前有一只大\VUgras{蛞蝓} \textsuperscript{(8)} 沉沉地爬着，
它刚吃掉一颗这里满地都是的好吃的小\VUgras{野草莓} \textsuperscript{(11)}。

%%K t09-c2-03-1 p
3. --- 地上铺着\VUgras{林中的花}：\VUgras{报春花} \textsuperscript{(9)}、
\VUgras{长春花} \textsuperscript{(10)} 和许多我不认得的植物，都在一丛丛
\VUgras{草} \textsuperscript{(12)} 中间开着。

%%K t09-c2-04-1 p
4. --- 在这一带，\VUgras{松露}差不多同\VUgras{蘑菇} \textsuperscript{(14)} 一样
常见；一个护林人的\VUgras{小猪} \textsuperscript{(13)} 正贪心地拱着找，看能不能
找到几个。

%%K t09-tit-10 sub
\VUsaut{3.0mm}
\VUpk{10.2pt}{40}{偷猎。}
\VUsaut{3.0mm}

%%K t09-c2-05-1 p
5. --- 我看见了一场很有趣的戏的\VUgras{结尾}。多亏他那只\VUgras{短腿猎犬}
\textsuperscript{(15)} 的鼻子，\VUgras{护林人} \textsuperscript{(16)} 刚好在
\VUgras{偷猎者} \textsuperscript{(22)} 要带走打到的\VUgras{野物} \textsuperscript{(17)}
时把他逮住。偷猎者宁可丢下猎物也不肯被抓，就跑了；\VUgras{野兔}
\textsuperscript{(18)}、\VUgras{母鹿} \textsuperscript{(19)}、\VUgras{狍}
\textsuperscript{(20)} 和\VUgras{野猪} \textsuperscript{(21)} 都留在草地上，叫护林人
好不高兴。

%%K t09-c2-06-1 p
6. --- 这林子里树木之多叫人吃惊。\VUgras{山毛榉} \textsuperscript{(23)} 和
\VUgras{白桦} \textsuperscript{(26)}，出\VUgras{松脂} \textsuperscript{(27)} 的
\VUgras{松树}，\VUgras{橡树} \textsuperscript{(29)}，还有许多别的树，枝叶都交在
一起。

%%K t09-c2-06-2 p
攀在高树干上的\VUgras{常春藤} \textsuperscript{(24)} 给这片\VUgras{乔林}
\textsuperscript{(25)} 添上鲜亮的绿，同\VUgras{苔} \textsuperscript{(31)} 那淡而柔和
的绿调和得很好；\VUgras{绿啄木鸟} \textsuperscript{(30)} 就在苔里找虫子。

%%K t09-tit-11 sub
\VUsaut{3.0mm}
\VUpk{10.2pt}{40}{林中的景色。}
\VUsaut{3.0mm}

%%K t09-c2-07-1 p
7. --- 那些老橡树叫我着迷。它们粗大盘曲的\VUgras{根} \textsuperscript{(32)} 半露
在土外，给它们添了一副极好的强壮气派；我不明白\VUgras{林主} \textsuperscript{(35)}
怎么忍心把它们放倒，交给\VUgras{伐木人}的\VUgras{锯} \textsuperscript{(34)}——
交给那\VUgras{伐木人} \textsuperscript{(33)}。那些可怜的\VUgras{树干}
\textsuperscript{(36)}，剥了\VUgras{皮} \textsuperscript{(37)} 的，或者叫\VUgras{斧}
\textsuperscript{(38)} 劈开的——那是\VUgras{短工} \textsuperscript{(39)} 的斧——
看着叫我难过。

%%K t09-c2-08-1 p
8. --- 近旁一个老\VUgras{烧炭人} \textsuperscript{(41)} 用他的\VUgras{锹}堆起一
\VUgras{堆} \textsuperscript{(42)} 木炭，一个老太婆背走一\VUgras{捆} \textsuperscript{(40)}
\VUgras{枯柴}，是砍倒的树上的枝子。

%%K t09-tit-12 sub
\VUsaut{3.0mm}
\VUpk{10.2pt}{40}{打猎。}
\VUsaut{3.0mm}

%%K t09-c2-09-1 p
9. --- 我本想到\VUgras{护林人的屋子} \textsuperscript{(46)} 去歇一会儿，可是走到
那个小\VUgras{池} \textsuperscript{(43)} 边——池上有一只好看的\VUgras{天鹅}
\textsuperscript{(45)} 在游——我发现前不久的一场\VUgras{大水} \textsuperscript{(44)}
把路变成了一片\VUgras{沼泽}。我只好绕道；经过林边那棵\VUgras{雪松}
\textsuperscript{(49)}、那些\VUgras{白杨} \textsuperscript{(48)} 和那棵\VUgras{榆树}
\textsuperscript{(47)} 时，我看见一只极好的\VUgras{雄鹿} \textsuperscript{(50)} 从
\VUgras{密丛} \textsuperscript{(54)} 里冲出来，后面紧追着一群不肯放松的
\VUgras{猎狗} \textsuperscript{(51)}，是\VUgras{号角} \textsuperscript{(53)} 在催
它们，那号角是\VUgras{骑马的猎手} \textsuperscript{(52)} 吹的。那可怜的畜生已经力尽了。它在离我
几步的地方倒下断了气。

%%K t09-c2-10-1 p
10. --- 护林人的儿子给\VUgras{猎}的声音引了出来，我们两个又走进林子。他
看见一只\VUgras{喜鹊} \textsuperscript{(56)} 停在一棵老橡树的枝上。他想它的窝
准藏在\VUgras{枝叶} \textsuperscript{(58)} 里，要去掏。

%%K t09-c2-10-2 p
他像\VUgras{松鼠} \textsuperscript{(61)} 一样灵便，从这根\VUgras{枝} \textsuperscript{(55)}
爬到那根枝，可是什么也没找着，只惊起一只老\VUgras{乌鸦} \textsuperscript{(60)}，
它正停在一根\VUgras{大枝} \textsuperscript{(57)} 上。
\VUsaut{4.0mm}
\VUfilet{22mm}
